\documentclass[10pt]{article}

% Specify the margins and text size.
\setlength{\textwidth}{6.4in}
\setlength{\textheight}{9.5in}
\setlength{\oddsidemargin}{0pt}
\setlength{\evensidemargin}{0pt}
\setlength{\topmargin}{0pt}
\setlength{\hoffset}{.05in}
\setlength{\voffset}{-1in}

\setlength{\parskip}{5pt}
\setlength{\parindent}{0pt}

% Load some fonts and symbol packages
\usepackage{latexsym}
\usepackage{pifont}       % contains 'star' symbol for counterinsurgency handbook title
\usepackage{yfonts} 
\usepackage{amsmath}
\usepackage{amsfonts}

\usepackage{graphicx}     % actually, this is loaded by pstricks
\usepackage[T1]{fontenc}
\usepackage{ifthen}
\usepackage{pstricks,pst-grad,pst-text,pst-node,multido,pst-plot,calc,pst-3dplot}
%\usepackage[all]{xy}
%\usepackage{animate}

% The hyperref package inserts links into the pdf.
\definecolor{MyLinkColor}{rgb}{.1,.2,1}
\definecolor{MyCiteColor}{rgb}{.1,1,.2}
\definecolor{MyURLColor}{rgb}{.4,.4,.4}
\usepackage[backref=true,pagebackref=false,hyperindex,colorlinks=true,
  linkcolor=MyLinkColor,urlcolor=MyURLColor]{hyperref}


% The tweaklist package is something I found on the web.  It provides a simple interface
% for making changes to spacing used in the itemize and enumerate environments.  Comment
% this out if you don't care to use tweaklist.
\usepackage{tweaklist}
\renewcommand{\itemhook}{\setlength{\parskip}{2pt}\setlength{\parsep}%
{1pt}\setlength{\topsep}{0pt}\setlength{\itemsep}{0pt}}

\newcommand{\U}{\underline{\hspace{5pt}}}

\usepackage{listings}
\newcommand{\Z}{\hphantom{0}}

\begin{document}
\pagestyle{empty}
\lstset{language=R, showspaces=false, showstringspaces=false}
\href{http://www.su.edu}{\includegraphics[height=1.75cm]{sulogo.eps}}
\vspace{-1.69cm}

{{\ }\hfill\small
\begin{tabular}{cl}
& Math 207\\
& Introduction to Statistics\\
\end{tabular}
}
\setlength{\baselineskip}{1.05\baselineskip}

\begin{center}
\textbf{\large  Quiz:  Central Limit Theorem}
\end{center}

\newcommand{\zobox}[2]{\begin{pspicture}(-0.2,0.2)(2.4,1.3)
\rput[r](-0.08,0.35){\small #1\vphantom{,'}}\psframe(0.0,0.05)(0.5,0.65)\rput(0.25,0.35){0}
   \rput[l](0.58,0.35){s\vphantom{',1}}
\rput[r](1.74,0.35){\small #2\vphantom{,'}}\psframe(1.8,0.05)(2.3,0.65)\rput(2.05,0.35){1}
   \rput[l](2.38,0.35){s\vphantom{',1}}
%\psframe(1,0.05)(1.6,0.65)\rput(1.3,0.35){1}\rput[bl](1.64,0.64){\scriptsize #2}
\psline(-0.9,0.9)(-0.9,-0.2)(2.8,-0.2)(2.8,0.9)(-0.9,0.9)
\end{pspicture}}
\newcommand{\threebox}[3]{\begin{pspicture}(0.2,0.2)(2.5,1.1)
\psframe(0,0)(0.7,0.7)\rput(0.35,0.35){#1}
\psframe(1,0)(1.7,0.7)\rput(1.35,0.35){#2}
\psframe(2,0)(2.7,0.7)\rput(2.35,0.35){#3}
\psline(-0.3,0.9)(-0.3,-0.2)(3,-0.2)(3,0.9)(-0.3,0.9)
\end{pspicture}}

\newcommand{\Zc}{\hphantom{0,}}



1. Complete the following table for the coin-tossing game.\vspace{-15pt}
\begin{center}
{\setlength{\tabcolsep}{15pt}\begin{tabular}{cc|cc|ccc}
              &\hbox{\hspace{.3in}}& \ \hspace{.5in}Number of heads\span &\hbox{\hspace{.1in}}
    & \ Percent of heads\span\\[5pt]
Number of     && Expected       &     && Expected & \\
Tosses        && Value          & SE  && Value & SE\\[3pt]\hline
&& & && &\\
\Zc\Z\Z\Zc100 && 50 & 5 && 50\% & 5\%\vphantom{\LARGE Y}\\[.2in]
\Zc\Z\Z2,500 &&    &   &&      &\\[.2in]
\Zc\Z10,000  &&    &   &&      &\\[.2in]
1,000,000    &&    &   &&      &\\[.2in]
\end{tabular}}
\end{center}
\vspace{1.2in}

2. Would taking the average of 25 measurements divide the likely size of the chance error by a factor of 5, 10 or 25? Justify your answer.
\vspace{1in}

3. A population of 1000 has 470 voters for candidate $A$ and 530 voters for candidate $B$.
A newspaper plans to poll a certain number of voters to try to estimate the percentage of 
voters for $A$.  Here is a box model.\vspace{-20pt}

\begin{center}
\zobox{530}{470}
\end{center}\vspace{-5pt}

\hspace{20pt} a) Compute $\mbox{EV}_{\scriptsize\%}$ and $\mbox{SD}_{\scriptsize\%}$
\vspace{1.5in}

\hspace{20pt} b) 
Compute 
$\mbox{SE}_{\scriptsize\%}$ for $n=16$, $n=25$, $n=49$ and $n=100$.  \vspace{-3pt}
%\hspace{20pt}\hphantom{a) } Note:  $1/8\approx 0.13$ and $1/14\approx 0.07$.
\vfill
\eject
{\ }


4. Nine hundred draws will be made at random with replacement
from the box shown below.
\begin{center}
\threebox{0}{0}{1}
\end{center}
\medskip

\hspace{10pt} a) 
The number of 1s among the draws will be around \underline{\hspace{45pt}} 
give or take \underline{\hspace{45pt}}  or so.
\vspace{1.25in}

\hspace{10pt} b) The chance of getting between 280 and 320 1s is \underline{\hspace{45pt}}.
\vspace{1.25in}

\hspace{10pt} c) What is the chance that the number of 1s will be less than 270?
\vspace{1.1in}

\hspace{10pt} d) 
The percentage of 1s among the draws will be around \underline{\hspace{45pt}} 
give or take \underline{\hspace{45pt}}  or so.
\vspace{1.25in}

\hspace{10pt} e) The chance of getting between 30\% and 36\% 1s is \underline{\hspace{45pt}}.
\vspace{1.25in}


\hspace{10pt} f) What is the probability that the percentage of 1s will 
be over 35\%.



\vfill
\eject



\end{document}

1. A population of 1000 has 470 voters for candidate $A$ and 530 voters for candidate $B$.
A newspaper plans to poll a certain number of voters to try to estimate the percentage of 
voters for $B$.  Here is a box model.

\begin{center}
\zobox{530}{470}
\end{center}

\hspace{20pt} a) Compute $\mbox{EV}_{\scriptsize\%}$ for the box.
\vspace{.5in}

\hspace{20pt} b) Given that $\mbox{SD}_{\scriptsize box}\approx 0.5$, compute 
$\mbox{EV}_{\scriptsize\%}$ for $n=16$, $n=25$, $n=49$ and $n=100$.  \vspace{-3pt}

\hspace{20pt}\hphantom{a) } Note:  $1/8\approx 0.13$ and $1/14\approx 0.07$.
\vspace{1.5in}


\vfill
\eject

\end{document}

